\chapter{Introduction}
\label{ch:introduction}

% Chapter 1: Introduction
% Typical structure: background, motivation, contributions, outline

\section{Background}
\label{sec:background}

Write the background and context of your research here. Introduce the broader
research area and explain why it matters.

For example, you might cite a foundational work like this~\cite{vaswani2017attention}.

Table~\ref{tab:example} shows an example of experimental results.

\begin{table}[h]
  \centering
  \caption{Example experimental results on benchmark datasets.}
  \label{tab:example}
  \begin{tabular}{lccc}
    \toprule
    Method & Accuracy & Params (M) & FLOPs (G) \\
    \midrule
    Baseline & 76.1 & 25.6 & 4.1 \\
    Ours     & \textbf{78.3} & 22.1 & 3.8 \\
    \bottomrule
  \end{tabular}
\end{table}

Figure~\ref{fig:example} illustrates the overall architecture.

\begin{figure}[h]
  \centering
  \fbox{\parbox{0.7\textwidth}{\centering\vspace{2cm}Architecture Diagram Placeholder\vspace{2cm}}}
  \caption{Overview of the proposed architecture.}
  \label{fig:example}
\end{figure}

\section{Motivation}
\label{sec:motivation}

Describe the specific problem you are addressing and why existing solutions
are insufficient.

\section{Contributions}
\label{sec:contributions}

Summarize the main contributions of your thesis:
\begin{enumerate}
  \item First contribution description.
  \item Second contribution description.
  \item Third contribution description.
\end{enumerate}

\section{Thesis Outline}
\label{sec:outline}

The remainder of this thesis is organized as follows:
\begin{itemize}
  \item \textbf{Chapter~\ref{ch:introduction}} introduces the research background and motivation.
  % \item \textbf{Chapter~\ref{ch:related}} reviews related work.
  % \item \textbf{Chapter~\ref{ch:method}} describes the proposed method.
  % \item \textbf{Chapter~\ref{ch:experiments}} presents experimental results.
  \item \textbf{Chapter~\ref{ch:conclusion}} concludes the thesis with a discussion of future work.
\end{itemize}
